\documentclass[11pt]{article}

% =============================================================
% CS229 Modern Extensions, Problem Set M1 (independent edition)
% This document is not an official Stanford assignment.
% =============================================================
\usepackage[letterpaper,margin=0.90in]{geometry}
\usepackage[T1]{fontenc}
\usepackage{lmodern}
\usepackage{amsmath,amssymb,amsthm,mathtools,bm}
\usepackage{enumitem}
\usepackage{fancyhdr}
\usepackage{hyperref}
\usepackage{xcolor}
\usepackage{microtype}

\hypersetup{
  colorlinks=true,
  urlcolor=blue!55!black,
  linkcolor=black,
  pdftitle={CS229 Modern Extensions: Problem Set M1},
  pdfauthor={Independent Study Edition}
}

\newcommand{\R}{\mathbb{R}}
\newcommand{\E}{\mathbb{E}}
\newcommand{\N}{\mathcal{N}}
\newcommand{\KL}{D_{\mathrm{KL}}}
\newcommand{\Cov}{\operatorname{Cov}}
\newcommand{\Law}{\operatorname{Law}}
\newcommand{\Normal}{\mathcal{N}}
\newcommand{\dd}{\,\mathrm{d}}
\newcommand{\given}{\,\middle|\,}
\newcommand{\hint}[1]{\par\vspace{2pt}\noindent\textit{Hint.} #1\par}
\newcommand{\pts}[1]{\hfill\textnormal{[\,#1 points\,]}}

\setlist[enumerate,1]{label=\textbf{(\alph*)},leftmargin=2.0em,itemsep=0.70em,topsep=0.35em,parsep=0.25em}
\setlength{\parskip}{0.35em}
\setlength{\parindent}{0pt}
\allowdisplaybreaks[1]
\pagestyle{fancy}
\fancyhf{}
\fancyhead[L]{\small CS229 Modern Extensions}
\fancyhead[R]{\small Problem Set M1}
\fancyfoot[C]{\thepage}
\renewcommand{\headrulewidth}{0.4pt}
\setlength{\headheight}{14pt}

\begin{document}
\begin{center}
  {\Large\bfseries CS229: Modern Machine Learning Extensions}\\[0.35em]
  {\LARGE\bfseries Problem Set M1}\\[0.2em]
  {\large Denoising Diffusion Probabilistic Models}\\[0.55em]
  {\small Independent Study Edition \textbullet{} Autumn 2026 \textbullet{} 100 points}
\end{center}
\vspace{0.2em}
\hrule
\vspace{0.75em}

\textbf{Scope.} This assignment studies Gaussian forward diffusion, exact reverse conditionals,
variational training objectives, noise-prediction parameterizations, and the implementation of a
small denoising diffusion probabilistic model (DDPM). It is aligned primarily with the diffusion
material in the \href{https://cs229.stanford.edu/main_notes.pdf}{2026 CS229 lecture notes}
(Chapter 14), with the \href{https://cs229.stanford.edu/index.html-spr26}{Spring 2026 course}
and \href{https://github.com/az9713/cs229-machine-learning-notes}{independent supplementary notes}
as additional context. This problem set was independently prepared and is \emph{not} an official
Stanford assignment.

\textbf{Instructions.} There are five problems, worth 100 points in total. Give mathematically
justified answers to all theoretical questions. Standard facts about multivariate Gaussian
distributions, Jensen's inequality, and KL divergence may be used without proof unless otherwise
stated. For programming problems, implement the indicated mathematical procedures rather than
using an existing diffusion-model training or sampling pipeline.

\textbf{Programming materials.} The dataset, numerical hyperparameters, starter code, and detailed
experimental specifications for Problems 4 and 5 will be supplied separately. The present version
contains the complete mathematical questions and the programming-task statements only.

\textbf{Notation.} For any \(d\ge 1\), \(\N(x;\mu,\Sigma)\) denotes the Gaussian \emph{density}
evaluated at \(x\), while \(\N(\mu,\Sigma)\) denotes the associated probability distribution.
The notation \(q\) is reserved for the fixed forward process; \(p_\theta\) denotes a learned
generative process. Empty products are equal to 1. All random vectors are column vectors.

\newpage
\section*{Problem 1. Gaussian Forward Diffusion \pts{15}}

Let \(X_0\sim q_{\mathrm{data}}\) be an arbitrary random vector in \(\R^d\). Define the
forward Markov chain \(X_0,X_1,\ldots,X_T\) by
\begin{equation}\label{eq:forward}
  q(x_t\mid x_{t-1})
  =\N\!\left(x_t;\sqrt{1-\beta_t}\,x_{t-1},\beta_t I_d\right),
  \qquad t=1,\ldots,T,
\end{equation}
where \(0<\beta_t<1\). Write
\[
  \alpha_t=1-\beta_t,\qquad
  \bar\alpha_t=\prod_{s=1}^t\alpha_s,\qquad
  \bar\alpha_0=1.
\]

\begin{enumerate}
  \item \pts{7} Prove, for every \(t\in\{1,\ldots,T\}\), that
  \begin{equation}\label{eq:forward-marginal}
    q(x_t\mid x_0)
    =\N\!\left(x_t;\sqrt{\bar\alpha_t}\,x_0,
                      (1-\bar\alpha_t)I_d\right).
  \end{equation}
  Deduce a method for sampling \(X_t\) conditional on \(X_0=x_0\) that does not require
  simulating the intermediate states \(X_1,\ldots,X_{t-1}\).

  \hint{Consider writing \(X_t\) as an affine function of \(X_0\) and independent standard
  Gaussian vectors. A telescoping identity simplifies the resulting covariance.}

  \item \pts{4} Assume \(X_0\) has mean \(m_0\) and covariance \(\Sigma_0\).
  Find \(\E[X_t]\) and \(\Cov(X_t)\). Does the Gaussian form of
  \(q(x_t\mid x_0)\) imply that the \emph{unconditional} distribution of \(X_t\) is
  Gaussian? Prove your claim or give a counterexample.

  \item \pts{4} Now extend the schedule to all positive integers. Give a necessary and
  sufficient condition on \(\{\beta_t\}_{t\ge1}\) for \(\bar\alpha_T\to0\).
  Prove that, under this condition, for each fixed \(x_0\in\R^d\),
  \[
    \Law(X_T\mid X_0=x_0)\ \xrightarrow[T\to\infty]{\mathrm{d}}\ \N(0,I_d).
  \]
  Is the assumption \(\beta_t>0\) for all \(t\) sufficient? Justify your answer.

  \hint{Take logarithms of the product defining \(\bar\alpha_T\).}
\end{enumerate}

\newpage
\section*{Problem 2. Reverse Conditionals and Variational Inference \pts{25}}

Consider a reverse-time generative model
\begin{equation}\label{eq:reverse-joint}
  p_\theta(x_{0:T})
  =p(x_T)\prod_{t=1}^{T}p_\theta(x_{t-1}\mid x_t),
  \qquad p(x_T)=\N(x_T;0,I_d),
\end{equation}
where \(x_{0:T}=(x_0,\ldots,x_T)\) and
\[
  p_\theta(x_{t-1}\mid x_t)
  =\N\!\left(x_{t-1};\mu_\theta(x_t,t),\sigma_t^2 I_d\right),
  \qquad \sigma_t^2>0.
\]
The objective is to maximize the marginal log-likelihood \(\log p_\theta(x_0)\)
for clean observations \(x_0\sim q_{\mathrm{data}}\).

\begin{enumerate}
  \item \pts{9} For \(2\le t\le T\), derive the conditional distribution
  \(q(x_{t-1}\mid x_t,x_0)\). Show that it is Gaussian, and obtain explicit formulas for
  \(\widetilde\mu_t(x_t,x_0)\) and \(\widetilde\beta_t\) such that
  \begin{equation}\label{eq:posterior}
    q(x_{t-1}\mid x_t,x_0)
    =\N\!\left(x_{t-1};\widetilde\mu_t(x_t,x_0),
                  \widetilde\beta_t I_d\right).
  \end{equation}
  Explain why this conditional distribution can be derived analytically even when the
  reverse marginal \(q(x_{t-1}\mid x_t)\) is not available in closed form.

  \hint{Use the Markov property and the proportionality
  \(q(x_{t-1}\mid x_t,x_0)\propto q(x_t\mid x_{t-1})q(x_{t-1}\mid x_0)\),
  then complete the square.}

  \item \pts{10} Starting from \(\log p_\theta(x_0)\), construct a variational lower bound
  \(\mathcal L_{\mathrm{ELBO}}(x_0;\theta)\) using
  \(q(x_{1:T}\mid x_0)\) as the variational distribution. Prove the decomposition
  \begin{align}\label{eq:elbo-decomposition}
    -\mathcal L_{\mathrm{ELBO}}(x_0;\theta)
    ={}&\KL\!\left(q(x_T\mid x_0)\,\|\,p(x_T)\right)\\
    &+\sum_{t=2}^{T}\E_{q(x_t\mid x_0)}\!\left[
      \KL\!\left(
        q(x_{t-1}\mid x_t,x_0)\,\|\,
        p_\theta(x_{t-1}\mid x_t)
      \right)
    \right]\notag\\
    &-\E_{q(x_1\mid x_0)}\!\left[
       \log p_\theta(x_0\mid x_1)
    \right].\notag
  \end{align}
  State which conditional-independence assumptions justify the factorization used in your
  derivation. Explain the probabilistic role of each term.

  \hint{Factor \(q(x_{1:T}\mid x_0)\) in reverse chronological order, beginning with
  \(q(x_T\mid x_0)\).}

  \item \pts{4} Evaluate the terminal term
  \[
    \KL\!\left(q(x_T\mid x_0)\,\|\,\N(0,I_d)\right)
  \]
  in closed form for fixed \(x_0\). Determine its limit as \(\bar\alpha_T\to0\).
  Explain why, with a fixed noise schedule and prior, it can be omitted from
  optimization over \(\theta\) without assuming that it is zero.

  \item \pts{2} Why does the posterior formula in part (a) become degenerate if one
  formally substitutes \(t=1\)? If
  \(p_\theta(x_0\mid x_1)=\N(x_0;\mu_\theta(x_1,1),\sigma_1^2 I_d)\),
  with \(\sigma_1^2>0\), express the reconstruction contribution as an expected
  negative log-likelihood and explain why it is treated separately.
\end{enumerate}

\newpage
\section*{Problem 3. Noise Prediction and the DDPM Objective \pts{20}}

By Problem 1, the forward process admits the representation
\begin{equation}\label{eq:reparam}
  X_t=\sqrt{\bar\alpha_t}\,X_0
      +\sqrt{1-\bar\alpha_t}\,\varepsilon,
  \qquad \varepsilon\sim\N(0,I_d),\quad \varepsilon\perp X_0.
\end{equation}
Let \(\varepsilon_\theta:\R^d\times\{1,\ldots,T\}\to\R^d\)
be a timestep-conditioned neural network.

\begin{enumerate}
  \item \pts{6} Rewrite the posterior mean \(\widetilde\mu_t(x_t,x_0)\), for \(t\ge2\),
  in terms of \(x_t\), the Gaussian noise \(\varepsilon\) in \eqref{eq:reparam},
  and the schedule parameters. Hence derive a parameterization of
  \(\mu_\theta(x_t,t)\) in terms of \(\varepsilon_\theta(x_t,t)\).

  \hint{Solve \eqref{eq:reparam} for \(X_0\), then substitute into the expression
  from Problem 2(a).}

  \item \pts{7} For \(t\ge2\), assume that the model variance is fixed at
  \(\sigma_t^2=\widetilde\beta_t\). Show that the corresponding KL term in
  \eqref{eq:elbo-decomposition} has the form
  \begin{equation}\label{eq:weighted-loss}
     \mathcal L_{t-1}(\theta)
     =w_t\,\E_{X_0,\varepsilon}\!\left[
       \left\|\varepsilon-
       \varepsilon_\theta(X_t,t)\right\|_2^2
     \right],
  \end{equation}
  and determine \(w_t\) explicitly. Here \(X_t\) is generated using
  \eqref{eq:reparam} for the stated timestep.

  The simplified noise-prediction objective is
  \begin{equation}\label{eq:simple-loss}
    \mathcal L_{\mathrm{simple}}(\theta)
    =\E_{\substack{X_0\sim q_{\mathrm{data}},\\
             t\sim\operatorname{Unif}\{1,\ldots,T\},\\
             \varepsilon\sim\N(0,I_d)}}\!\left[
      \left\|\varepsilon-
      \varepsilon_\theta(X_t,t)\right\|_2^2
    \right].
  \end{equation}
  Explain why removing the timestep-dependent weights generally changes the optimization
  problem when the same network parameters are shared across timesteps. Give a simple
  mathematical example illustrating this difference. In your discussion, distinguish the
  \(t\ge2\) KL terms from the \(t=1\) reconstruction term.

  \item \pts{5} Fix a timestep \(t\), and suppose \(f\) ranges over all measurable
  functions with finite expected squared error. Find the minimizer (up to almost-sure
  equality) of
  \[
      \min_f\ \E\!\left[\|\varepsilon-f(X_t)\|_2^2\right].
  \]
  Prove your answer. Explain why minimizing this loss does not, in general, recover the
  particular noise realization used to produce a given \(X_t\).

  \hint{Decompose \(\varepsilon-f(X_t)\) around
  \(\E[\varepsilon\mid X_t]\) and use orthogonality of conditional-expectation residuals.}

  \item \pts{2} Describe how the learned noise predictor can be used to sample
  approximately from the data distribution, beginning with
  \(X_T\sim\N(0,I_d)\). State the reverse transition used for \(t\ge2\).
  For the final step \(t=1\), explain the distinction between sampling a
  positive-variance Gaussian decoder and using a deterministic reverse-mean update
  (a common implementation convention).
\end{enumerate}

\newpage
\section*{Problem 4. Numerical Verification of the Diffusion Process \pts{15}}

In this problem you will implement the analytic constructions from Problems 1 and 2,
and validate them numerically. Use NumPy for the required computations.
\emph{Detailed input data, numerical schedules, and starter-code interfaces will be
released separately.}

\begin{enumerate}
  \item \pts{3} Implement a function that, given \(\beta_1,\ldots,\beta_T\),
  constructs \(\alpha_t\) and \(\bar\alpha_t\). Document the indexing convention
  and validate that the supplied values satisfy \(0<\beta_t<1\).

  \item \pts{7} Implement (i) direct sampling from \(q(x_t\mid x_0)\) without
  simulating the full chain and (ii) the analytical mean and covariance of
  \(q(x_{t-1}\mid x_t,x_0)\) for \(t\ge2\). The functions should support
  batches with different timesteps for different observations. Explain how you
  handle broadcasting and tensor shapes.

  \item \pts{5} Design and perform Monte Carlo tests comparing empirical and
  theoretical conditional means and covariance matrices at several timesteps.
  Independently compare direct sampling with sampling obtained by composing
  one-step forward transitions. Provide an appropriate numerical error measure,
  discuss sampling variability, and explain what the tests do and do not establish.
\end{enumerate}

\newpage
\section*{Problem 5. Implementing a DDPM from Scratch \pts{25}}

You will now implement and train a denoising diffusion probabilistic model on a
low-dimensional, multimodal synthetic distribution. The dataset specification,
baseline hyperparameters, and evaluation protocol will be supplied separately.
Use PyTorch for this problem. Do not use a ready-made DDPM trainer or sampler.

\begin{enumerate}
  \item \pts{5} Design a timestep-conditioned neural network
  \(\varepsilon_\theta(x_t,t)\) that predicts the noise used in the forward
  corruption process. Describe the choice of timestep representation and explain
  why time conditioning is required.

  \item \pts{5} Implement a stochastic estimator of the simplified objective
  \eqref{eq:simple-loss}. Your implementation should sample timesteps and Gaussian
  perturbations, construct noisy inputs directly, and use automatic differentiation
  to update \(\theta\). Explain why the sampled minibatch loss is an estimator
  of the stated population objective.

  \item \pts{7} Implement a reverse-diffusion sampler, starting from
  \(X_T\sim\N(0,I_d)\). Generate observations from the trained model and
  compare their distribution with held-out samples from the target distribution.
  Visualize intermediate stages of the reverse process and interpret the result
  in terms of the learned reverse transition kernels.

  \item \pts{8} Carry out a controlled comparison between two distinct noise
  schedules with the same terminal value \(\bar\alpha_T\). Derive any schedule
  parameter needed to enforce this equality. Evaluate denoising accuracy and
  sample quality using quantitative metrics appropriate to the provided dataset.
  Discuss whether lower noise-prediction error necessarily implies better
  generative performance, and identify limitations of the experimental comparison.
\end{enumerate}

\vspace{0.6em}
\hrule
\vspace{0.6em}
{\small
\textbf{Suggested references.}
\href{https://cs229.stanford.edu/main_notes.pdf}{Ma and Ng, \emph{CS229 Lecture Notes} (2026), Chapter 14};
\href{https://arxiv.org/abs/2006.11239}{Ho, Jain, and Abbeel, \emph{Denoising Diffusion Probabilistic Models} (2020)};
\href{https://github.com/az9713/cs229-machine-learning-notes}{independent Spring 2026 notes}.
}

\end{document}
